\documentclass[10pt,a4paper]{article} \usepackage{fontspec} \usepackage{polyglossia} \setmainlanguage{ukrainian} \setotherlanguage{english} \setmainfont{DejaVu Sans} \newfontfamily\cyrillicfont{DejaVu Sans} \usepackage[left=2cm,right=2cm,top=2cm,bottom=2cm]{geometry} \usepackage{enumitem} \usepackage{titlesec} \usepackage{hyperref} \usepackage{parskip} \hypersetup{ colorlinks=true, linkcolor=black, urlcolor=blue, pdfauthor={Anna Sheludko}, pdftitle={Анна Шелудько — Резюме} } \titleformat{\section} {\normalfont\large\bfseries} {}{0em}{} [\vspace{0.2em}\hrule height 0.6pt] \titlespacing*{\section}{0pt}{1em}{0.5em} \pagestyle{empty} \setlist[itemize]{ left=1.2em, nosep, topsep=3pt, itemsep=2pt } \begin{document} \begin{center} {\LARGE\bfseries Анна Шелудько}\\[0.35em] Data Engineer \textbar{} Python Developer\\[0.3em] annasheludko152@gmail.com\\ \href{https://anna1sheludko.github.io}{anna1sheludko.github.io} \textbar{} \href{https://github.com/anna1sheludko}{github.com/anna1sheludko}\\[0.3em] Українка \textbar{} проживаю в Польщі \textbar{} відкрита до співпраці \end{center} \section{Профіль} Інженерка даних та Python-розробниця, яка спеціалізується на проєктуванні ETL-процесів, роботі з реляційними базами даних та автоматизації обробки даних. Створюю рішення, що перетворюють розрізнені джерела даних на структуровані та готові до аналізу бази даних. Працюю з Python, SQL, PostgreSQL і Docker. Нещодавно реалізувала ETL-пайплайн для обробки 1,5 мільйона записів у 9 таблицях із валідацією схеми, автоматичними тестами та розгортанням у Docker. \section{Технічні навички} \begin{itemize} \item \textbf{Мови програмування:} Python, SQL, PL/SQL, Java, C++, R \item \textbf{Бази даних:} PostgreSQL, Oracle, MySQL, MS SQL Server, NoSQL, семантичні бази даних, графи знань \item \textbf{Обробка даних:} Pandas, NumPy, Pandera, Excel, CSV, JSON \item \textbf{Інфраструктура та DevOps:} Docker, Git, GitHub Actions, Linux, YAML \item \textbf{Тестування та якість:} pytest, валідація даних, тестування ETL-процесів, забезпечення якості даних \item \textbf{Також маю досвід з:} REST API, Apache Airflow (базовий рівень), статистикою, машинним навчанням, NLP, комп'ютерними мережами та хмарними технологіями \end{itemize} \section{Проєкти} \textbf{End-to-End ETL Pipeline} \hfill \href{https://github.com/anna1sheludko/end-to-end-pipeline}{GitHub} \vspace{0.3em} \textit{Python, Pandas, PostgreSQL, Docker, Pandera, pytest} \begin{itemize} \item Розробила та впровадила повноцінний ETL-пайплайн для даних електронної комерції, що охоплює 9 таблиць та близько 1,5 мільйона записів. \item Реалізувала механізм автоматичних повторних спроб під час отримання даних, валідацію схеми за допомогою Pandera та високопродуктивне завантаження через PostgreSQL COPY, що забезпечило приблизно 8-кратне пришвидшення порівняно зі стандартними INSERT. \item Налаштувала конфігурацію в YAML, структуроване логування з ротацією файлів, автоматизовані тести та CI/CD через GitHub Actions. \item Проєкт повністю відтворюється за допомогою Docker Compose. \end{itemize} \vspace{0.6em} \textbf{API $\rightarrow$ PostgreSQL Pipeline} \hfill \textit{у процесі розробки} \vspace{0.3em} \textit{Python, REST API, PostgreSQL, Docker} \begin{itemize} \item Розробляю ETL-пайплайн, який отримує дані з публічного REST API, виконує їхню валідацію та завантажує до PostgreSQL для подальшого аналізу. \item Проєкт включає автоматизацію збору даних, контроль їхньої якості та контейнеризацію за допомогою Docker. \end{itemize} \section{Освіта} \textbf{Інженерія даних (бакалаврат)} \hfill \textit{очікуване завершення: 2026} \vspace{0.2em} Свентокшиська політехніка, Кельце, Польща \begin{itemize} \item Ключові дисципліни: бази даних, сховища даних, семантичні бази даних, графи знань, машинне навчання, статистика, Python, Java, C++, комп'ютерні мережі та управління проєктами. \item Шестимісячне професійне стажування в межах навчальної програми. \end{itemize} \section{Мови} \textbf{Українська} (рідна) \textbar{} \textbf{Польська} (B2) \textbar{} \textbf{Англійська} (B2) \section{Доступність} Відкрита до співпраці над фріланс-проєктами, стажуваннями та junior-позиціями у сфері Data Engineering і Python Development. Надаю перевагу контакту електронною поштою та зазвичай відповідаю протягом 24 годин. \end{document}